\begin{tikzpicture}[
  font=\sffamily,
  line width=0.7pt,
  usecase/.style={
    draw,
    ellipse,
    align=center,
    text width=3.45cm,
    minimum height=1.1cm,
    inner sep=4pt,
    fill=white
  },
  actorlink/.style={draw}
]
  % Use-case diagram frame and heading.
  \draw (0,0) rectangle (17,10);
  \node[anchor=west,font=\bfseries] at (0.25,9.65)
    {uc [block] Enterprise Use Cases [EPS Development Enterprise]};
  \draw (0,9.3) -- (17,9.3);

  % Enterprise subject boundary.
  \draw (2.2,0.7) rectangle (13.1,8.75);
  \node[anchor=north,align=center,text width=9.7cm] at (7.65,8.48)
    {\textbf{Academic Nanosatellite EPS Development Enterprise}};

  % Mission-level enterprise use cases.
  \node[usecase] (maintain) at (4.8,6.85)
    {\textbf{Maintain EPS Evidence Chain}};
  \node[usecase] (configure) at (10.5,6.85)
    {\textbf{Configure Framework\\ for an EPS Case}};
  \node[usecase] (assess) at (4.8,4.75)
    {\textbf{Assess a Design Change}};
  \node[usecase] (reproduce) at (10.5,4.75)
    {\textbf{Reproduce an Analysis}};
  \node[usecase] (compare) at (4.8,2.55)
    {\textbf{Compare Candidate\\ EPS Architectures}};
  \node[usecase] (review) at (10.5,2.55)
    {\textbf{Review an Early\\ Design Justification}};

  % External reviewer actor. Its associations appear only where the actor
  % directly participates in the enterprise use case.
  \draw (15.05,5.55) circle (0.24);
  \draw (15.05,5.31) -- (15.05,4.55);
  \draw (14.62,5.05) -- (15.48,5.05);
  \draw (15.05,4.55) -- (14.62,4.02);
  \draw (15.05,4.55) -- (15.48,4.02);
  \node[align=center,text width=2.3cm] at (15.05,3.5)
    {\textbf{Reviewer}\\[-1pt]\footnotesize External actor};

  \draw[actorlink] (14.62,5.05) -- (reproduce.east);
  \draw[actorlink] (14.62,5.05) -- (review.east);
\end{tikzpicture}
